\Needspace{12\baselineskip}
\originalchapter{原子模型、超胞与缺陷能}
\originalsection{势函数如何成为物理假设}
3.320由Gerbrand Ceder和Nicola Marzari讲授原子尺度建模。它的主线从经验势、电子结构延伸到有限温度采样、动力学与粗粒化。本卷据公开讲义重编物理联系，不重复数学仓库的通用插值或线性代数章节。

给定原子坐标 $\mathbf R_i$，最简对势模型 $U=\sum_{i<j}v(r_{ij})$。Lennard--Jones势 $v(r)=4\epsilon[(\sigma/r)^{12}-(\sigma/r)^6]$ 以短程排斥和长程吸引描述某些非极性体系。参数适合一个物性不保证所有物性都准确；金属中的电子环境、共价键方向和电荷转移通常不能由同一个简单中心对势充分描述。

对势最低点 $r_0=2^{1/6}\sigma$，但晶体的平衡最近邻距离不必等于这个值，因为一个原子同时感受多层邻居。若FCC晶格常数为 $a$，最近邻距离 $a/\sqrt2$，每原子能量包含全晶格和而不是12个最近邻之外全删。截断也影响最低点，需要检查范围收敛。

嵌入原子法写 $U=\sum_iF(\rho_i)+\tfrac12\sum_{i\ne j}\phi(r_{ij})$，$\rho_i=\sum_{j\ne i}f(r_{ij})$。非线性嵌入项允许同一对原子的有效作用随环境变化。它不是把LJ对势换一个更长的公式，而是改变模型中可表达的多体信息；具体参数的训练范围仍限制缺陷、表面及不同相的可靠性。

\originalsection{周期边界和弛豫}
周期超胞表示无限周期排列的一个重复单元。正交盒的最小镜像位移 $\Delta\mathbf r=\mathbf r_i-\mathbf r_j-L\,\mathrm{round}[(\mathbf r_i-\mathbf r_j)/L]$，截断半径必须小于 $L/2$ 才能使用单镜像求和。非正交盒需按晶格向量处理，分量简单舍入并不总给最近镜像。

固定胞形优化内部坐标，与同时优化晶格大小和形状是不同约束。数值最小化通常只能找到某个局部极小，不证明全局最低结构。先检查力与能量导数一致，再讨论优化收敛；用错符号的力也可能给出看似稳定的输出。

\originalsection{形成能和表面能}
对含 $N$ 个完美晶格位点的胞，移去一个原子并送回体相，空位形成能是 $E_v=E_{N-1}-(N-1)E_N/N$。必须保持势模型、参考状态、边界与计算精度一致。增大超胞可以减少缺陷周期镜像间相互作用；弛豫后形成能通常降低，但若另改变体积，比较的物理约束也改变。

表面薄片在周期边界下常产生两个表面，表面能 $\gamma_s=(E_{\rm slab}-N e_{\rm bulk})/(2A)$。收敛有两个独立方向：物质薄片厚度和真空间隙。将真空加厚不能修复薄片内部过薄的问题。非对称表面需要分别考虑两侧以及电偶极校正，不能不加说明继续套两个等价面的公式。

\originalsection{从能量求力、压力和弹性}
对中心对势，$\mathbf F_{ij}=-v'(r_{ij})\mathbf r_{ij}/r_{ij}$，其中 $\mathbf r_{ij}=\mathbf R_i-\mathbf R_j$，总力成对相反。LJ力径向大小为 $(24\epsilon/r)[2(\sigma/r)^{12}-(\sigma/r)^6]$；在短程为排斥。有限差分能量检验必须选择非平衡、非对称构型，否则恰巧零力不能检出符号错误。附带程序使用明确位移后的构型做导数核验。

施加小均匀应变 $\varepsilon$，能量展开 $E=E_0+V\sigma_{ij}\varepsilon_{ij}+VC_{ijkl}\varepsilon_{ij}\varepsilon_{kl}/2+\cdots$。应力是一阶导、弹性常数二阶导；若内部原子位置可弛豫，应在每个应变点先按相同标准弛豫，得到的与固定内部坐标“clamped”常数不同。拟合太小应变会受能量噪声影响，太大应变又含高阶项，因此要比较拟合窗口。

中心对势常对某些弹性常数施加限制，而金属和共价材料未必满足。比较多种独立物性（晶格常数、凝聚能、弹性、缺陷和表面）可判断势函数的表达能力；只在一个训练点准确，不能证明模型可迁移。

\originalsection{离子长程力与不同材料的势}
离子晶体的 $q_iq_j/(4\pi\epsilon_0r)$ 长程和不能直接按LJ最小镜像截断处理。Ewald方法引入Gaussian屏蔽，将短程实空间和光滑倒空间相加，再扣自作用；屏蔽参数只是计算分拆，收敛结果应与其选择无关。体系电中性、边界电势约定和胞形影响某些长程项，必须在能量比较中一致。

共价键偏好角度，可用键序或显式三体角项；金属电子密度环境适合嵌入原子类多体模型；极性体系还需电荷或极化。选模型先问要保留什么自由度与机制，再问程序快不快。将跨相、断键和高压行为外推到未训练区间，是物理建模误差，不是积分精度误差。

\originalsection{FCC晶格和的解析基准}
若 $v=A/r^{12}-B/r^6$，定义最近邻距 $r_{nn}$ 和无量纲壳和 $S_p=\sum_j(r_{nn}/r_j)^p$，每原子能量为 $e=AS_{12}/(2r_{nn}^{12})-BS_6/(2r_{nn}^6)$。驻点给 $r_{nn}^6=2AS_{12}/(BS_6)$，$a=\sqrt2r_{nn}$。这与孤立对 $r_0^6=2A/B$ 的差异完全由壳和产生。

使用公开GULP手册的Au裸LJ参数，附带独立壳枚举给两、四、七层结果 $a\simeq4.08187,4.07857,4.07805$ \AA{}。这个计算是真实运行的解析势基准，不是GULP或真实金样品的测量。更远吸引壳改变能量与平衡尺度，说明截断收敛需要实际比较。

\originalsection{缺陷参考化学势与有限尺寸}
一般形成能写成 $E_f=E_{
m defect}-E_{
m perfect}-\sum_\alpha\Delta N_\alpha\mu_\alpha$。对空位 $\Delta N=-1$，送回体相的 $\mu=e_{\rm bulk}$ 给前述式；对替位杂质，既移去母元素又加入新元素，两个化学势都需要指定。带电缺陷还有电子化学势与电荷补偿/有限尺寸修正，不能照搬中性对势空位公式。

周期镜像引起弹性或静电相互作用。检查更大胞时，应保持缺陷浓度定义和边界约束清楚，区分“单缺陷稀释极限”与有限浓度周期阵列。源码中的模型示例只处理它明确实现的中性势问题，没有输出未经计算的实验形成能。
